% ====================================================================
% 使用说明
% ====================================================================
% 本文档为开题报告正文模板，请将以下示例内容替换为您的实际开题报告内容。
% 
% 1. 章节标题使用方法：
%    \section{一级标题}
%    \subsection{二级标题}
%    \subsubsection{三级标题}
%
% 2. 参考文献引用方法：
%    使用 \cite{引用标签} 命令，例如：\cite{knuth1984texbook}
%    多个文献引用：\cite{knuth1984texbook,lamport1994latex}
%
% 3. 公式添加方法：
%    行内公式：$E=mc^2$
%    独立公式：
%    \begin{equation}
%        E = mc^2
%    \end{equation}
%
% 4. 列表添加方法：
%    无序列表：
%    \begin{itemize}
%        \item 第一项
%        \item 第二项
%    \end{itemize}
%    
%    有序列表：
%    \begin{enumerate}
%        \item 第一项
%        \item 第二项
%    \end{enumerate}
% ====================================================================

\section{论文选题的依据：}

随着工程结构向轻量化与高性能化方向发展，数值仿真对计算精度和效率的要求日益提升。曲梁结构在桥梁、拱门及大跨度曲面屋顶等领域应用广泛；另外，屈曲构件也可类似于曲梁构件，当发生弯曲时也会出现清剪切效应。Euler-Bernoulli梁理论基于平截面假设，在处理厚梁或需考虑剪切变形时难以满足精度需求。相比之下，Timoshenko曲梁模型通过独立描述弯曲与剪切效应，为此类结构提供精准建模基础，该模型同时考虑弯曲变形和剪切变形的影响\cite{li1990}。但在数值离散时，若传统有限元单元采用不匹配的形函数和完全积分方案，会产生虚假剪切变形，导致系统呈现过高剪切刚度，从而引发剪切自锁现象\cite{hughes1977, belytschko1985, kreja1988, bouclier2012, zhang2018, baier-saip2020}。此现象源于位移场与转角场的自由度不匹配，使虚假剪切应变能主导弯曲响应，导致模型过度刚化。该问题在屈曲分析中亦会导致结果与实际存在显著误差。因此，解决曲梁剪切自锁问题及其在屈曲分析中的应用成为本研究核心课题。


为消除剪切自锁现象，学者基于单一位移场框架提出减缩积分技术\cite{zienkiewicz1971, kreja1988, cheng1997}，其核心是通过降低整个刚度矩阵的阶数（包括弯曲项和剪切项），避免高阶假性约束引入，从而允许单元更灵活地模拟纯弯曲模式而不产生寄生剪切变形。在减缩积分基础上发展的选择性减缩积分则将弯曲和剪切项分开处理\cite{hughes1977, zhang2018, bouclier2012}：对弯曲项使用完整积分保证刚度精度，对剪切项使用减缩积分消除自锁现象。基于此方法，Stolarski和Belytschko\cite{belytschko1985}与 Melo 等人\cite{demelo1992}针对曲单元与弯管结构，通过局部降阶积分方案减少高阶应变项的积分点，规避复杂几何下的锁闭风险。然而选择性减缩积分可能导致单元过度刚化而精度下降，而均匀减缩积分虽精度较好却易引发矩阵奇异。为此，应力投影方法被提出並验证\cite{sun2025, belytschko1985, elguedj2007, wriggers2023, dvorakova2022}，该方法通过将应力场投影到低阶空间修改单元刚度矩阵，在保持精度的同时避免过度刚化。研究表明该技术与模态分解在消除寄生应力方面具有等价性，但稳定性更优。另一方面，Bouclier\cite{bouclier2012}与Cheng等人\cite{cheng1997}将缩减积分与高阶基函数技术耦合处理剪切能项，通过单点高斯积分结合气泡函数空间增强旋转场插值精度，同步减少膜应变和剪切应变能的积分点以消除锁定。但需注意，积分技术可能引入零能模态，需通过稳定化方法修正以维持解的稳定性和精度。

鉴于积分技术对网格取向敏感，学者提出了改进的单一位移场框架方法。高阶插值法透過提升階數或引入厚度方向二次變化的轉動自由度，使剪切應變自然包含必要項\cite{li1990, tai2016, koziey1994, dawe1974}，有效减轻锁定效应，在粗网格或复杂几何中表现优异，並可耦合减缩积分进一步避免假性约束。这种耦合不仅提升收敛率，还能在不引入额外变量条件下平衡弯曲与剪切项。基于相同原理，混合插值法通过建立完整应变-位移关系，包含翘曲剪应力项，引入内部节点增加自由度，缓解插值不一致导致的伪剪切约束，並使用静态凝聚消除内部节点自由度，保持元件仅端节点输出\cite{schnabl2007, wang2014}，实现类似选择性积分的效应，在单一位移场框架下精确控制刚度矩阵，避免寄生变形並维持高精度。形函数方法中，陈太聪等\cite{LXXB200906019}与田荣\cite{LXXB201901027}分别提出构造精确形函数和$C^1$连续形函数算法，使粗网格能准确捕捉屈曲载荷和失稳模态，其方法可同时用于強/弱形式离散且无需网格加密即可实现高阶收敛。然而，形函数精确形函数需一般解推导，适用范围窄

为突破单一位移场框架对未知函数的严格约束，学者转向引入第二变量扩展函数空间。其代表方法是基于混合变分原理的混合有限元方法，把应力作为独立变量，与位移场共同变分，使方程满足约束指数，並通过赫林格-赖斯纳变分原理静止条件消去应力参数，得到标准刚度矩阵；对于曲梁的近似轴线构建均衡应力使其能精确计算刚度，在模拟任意形状曲梁或圆环，应力预测准确，收敛快速\cite{malkus1978, saleeb1987,benedetti1989}。但因其限于线弹性静力模型。基于此Saritas 和Hale等\cite{saritas2009, hale2012}引入非弹性材料响应和无网格方法，对于非弹性，沿梁长多截面数值积分材料响应並且在无网格情况下，直接施加本质边界条件，混合处理剪切应变避免锁定。这个方法虽能处理非弹性，但计算效率低且缺乏高阶连续和动态守恒。见及于此，\cite{ choi2023, gross2019}用IGA实现更高连续性、降低阶数大幅节省计算，並引入守恒方案确保长时动态稳定，通过独立场局部形函数设计，避免剪切锁定，並把增强了应变和不连续场处理截面锁定。，该方法同时离散位移场与应力场两个基本未知量，通过独立近似应力场与位移场放宽主变量约束条件，从变分原理层面消除剪切自锁现象。Xiao与McCarthy \cite{xiao2003}进一步提出基于局部变分原理的改进方案，通过将传统位移-旋转变量对替换为位移-剪切应变对，重构势能泛函的数学形式，从根源解除导致剪切锁定的约束条件。然而传统混合有限元在空间离散化时存在局限性：等阶插值组合易引发数值不稳定。为此，混合有限元法通过修改伽辽金框架\cite{franca1988}，添加最小平方类扰动项增强稳定性，使任意连续插值组合（包括等阶插值）均能收敛，且无需引入全局额外变量——该技术通过单元内部计算实现，对低阶元素尤为有效。

由于混合有限元需引入额外变量导致系统矩阵规模扩大，显著增加计算复杂度与资源消耗，且在位移/应力场插值空间匹配、多场耦合处理方面要求严苛。作为替代方案，假定应变场法同样基于混合变分原理，其方法先假设应变为常值或线性；积分得到包含应变模态与刚体模态的位移表达式；通过节点位移求解广义常数，转换至全局笛卡尔座标，彻底避开相容应变造成的寄生约束，且收敛特性优于修正等参元；二节点与三节点元均简单高效\cite{choit1995, prathap1993, LXXB201302021}。但局限明显仅适用特定曲梁；难以处理任意复杂截面；本构模型受梁理论限制。为此，将成熟的固壳技术系统移植至梁结构，提出8节点纯位移三维固梁元。这是从经典梁/曲梁元向完整三维连续体梁元的重大转型，在沿梁长方向采用标准全积分，两个厚度方向采用减缩积分，横向剪切与膜锁定用假设应变场组合抑制，使得本构自由度大幅提升能直接使用三维非线性材料模型\cite{frischkorn2013}。但其方法为传统有限元，刚度矩阵较密需多节点逼近圆形/曲面截面。因此，进一步利用投影法的高阶连续性与精确几何，结合假设应变。在系点评估相容剪切与膜应变；采用线性/封闭形式外插得到假设应变。得到膜/剪切锁定完全消除；厚度相关锁定自然缓解；单元素截面下精度极高；计算极其高效\cite{patton2025, greco2017}，但其計算成本高且幾何逼近誤差大。


尽管现有研究已发展出减缩积分、假定应变场、混合方法、高阶插值及无网格方法等多种技术以抑制曲梁分析中的剪切自锁现象，这些方法仍存在一些关键问题未解决，其共性局限可归纳为以下三点:
\begin{enumerate}
\item 设计过程依赖经验性选择，減縮積分、選擇性減縮積分及應力投影方法的核心在於人工選擇積分點數量，對彎曲項與剪切項的差別處理高度依賴經驗判斷，易引发刚度矩阵奇异，需额外稳定化处理
\item 受限于基于单元的离散框架，单一位移场框架对未知函数的严格相容性约束，而混合有限元也需要严格的位移/应力插值空间匹配，从而无法从根本上实现约束比的精确可控。
\item 缺乏对“约束比”的显式误差分析，缺乏核心的系统理论指导，未能建立其与数值误差之间的定量关系
\end{enumerate}
上述方法的局限使多场耦合等工程面临挑战。为此，本文将以曲梁为研究对象，重点探讨其混合离散分析中消除剪切自锁的“最优约束比例”问题，旨在通过理论分析与数值实验明确不同离散策略下约束条件与自锁行为之间的内在关联，从而填补现有方法在理论指导层面的空白。在此基础上，本文将进一步建立一种“自锁约束比例可控的有限元与无网格混合离散分析方法”，通过有机结合有限元的结构化离散优势与无网格法的灵活近似特性，实现约束比例的可控调节，在保证计算精度的同时有效避免自锁现象，为工程实践提供一种兼具鲁棒性与适应性的新型数值解决方案。



\section{研究特色及创新之处}
\begin{enumerate}
\item 考虑剪切应变与膜应变自锁以及剪切应力与膜应力自锁。
\item 推出了剪切应变与膜应变自锁以及剪切应力与膜应力自锁系数的估计表达式从而得到最优约束比例。
\item 提出了具有最优约束比例、曲梁及屈曲分析的免自锁有限元无网格混合离散控制方法。
\end{enumerate}
\section{研究内容}

本文将针对曲梁剪切自锁问题，研究位移场和约束场的自由度数量比与LBB稳定性条件影响机理，並利用最优约束比和无网格法的离散便利性，从而建立曲梁最优约束比的有限元无网格混合分析方法，並应用于曲梁结构屈曲分析中。具体内容如下：

\begin{enumerate}
    \item 基于Timoshenko梁理论，通过特征值分析与模态分解，量化不同自由度配比下离散inf-sup常数的变化规律，揭示当约束场自由度过高时易引发压力/应力振荡，重点推导约束比与位移误差、应力收敛性间的显式关系，为约束比优化奠定严谨理论基础。
    \item 设计混合离散方案：位移场采用再生核无网格形函数离散，支持全域连续近似；约束场采用有限元离散，通过泛函分析，建立曲梁分析的最优约束比
    \item 将最优约束比例引入曲梁有限元无网格混合离散分析分法，建立免自锁的有限元无网格混合离散分析方法，並将其引入曲梁的屈曲分析，从而构造一种高精度稳定的曲梁屈曲分析方法
\end{enumerate}

\section{预期成果}
本论文的预期成果如下：
\begin{enumerate}
\item 曲梁自锁问题的LBB稳定性条件系数估计。
\item 最优约束比例免自锁曲梁屈曲分析的有限元无网格混合离散方法。
\end{enumerate}

\section{进度安排}

\begin{table}[H]
\centering
\renewcommand{\arraystretch}{1.2}
\begin{tblr}{
    width = 0.95\linewidth,
    colspec = {X[1.2,l] X[2.3,l] X[0.8,c]},  %：移除逗號
    row{1} = {font=\bfseries, c, },  % 添加表頭背景色
    row{2-Z} = {rowsep=8pt},
    row{6} = {font=\normalsize},
    cell{7}{1} = {c=3}{r},
    hline{1,2,8} = {1pt}, % 顶部、表头下、底部的线
    % hline{3,4,5,7} = {0.5pt, solid}, % 行之间的细线
    vlines = {0pt}, % 不显示垂直线
}
\textbf{時間} & \textbf{研究內容} & \textbf{進度情況} \\
2025.7-2025.12 & 验证LBB稳定性条件与自由度比例之间的关系 & \scalebox{2}{$\circ$} \\
2026.1-2026.2 & 进行Timoshenko曲梁的无网格有限元混合离散分析 & \\
2026.3-2026.9 & 曲梁结构免自锁在屈曲分析中的应用 &  \\
2026.10-2026.12 & 整理成果並完成论文初稿撰写 &\\
2027.1-2027.3 & 进一步修改完善论文，形成並提交论文终稿 &  \\
\scalebox{2}{$\bullet$} 已完成 \quad \scalebox{2}{$\circ$} 正在進行 & & \\
\end{tblr}

\vspace{0.5em}

\end{table}

% ====================================================================
% 注意：请在 references.bib 文件中添加对应的参考文献条目
% ====================================================================